\chapter{TINJAUAN PUSTAKA DAN DASAR TEORI}                

\section{Tinjauan Pustaka}

Berdasarkan penelitian yang dilakukan oleh Solovyev et al. implementasi 
CNN pada FPGA untuk realtime processing dapat dilakukan dengan 
melakukan konversi \emph{floating-point} menjadi \emph{fixed-point} supaya lebih cocok dijalankan 
dalam skala hardware \cite{Solovyev2019FixedPointCNN}. Penelitian ini 
menggunakan data MNIST dalam melakukan pelatihan model. Selama proses 
pelatihan model dilakukan augmentasi data berupa rotasi acak 10 derajat, inversi warna 
dan \emph{noise} acak dari $0\%$ sampai $10\%$. Hasil pelatihan parameter model 
memory disimpan sebagai blok memori. Hasil akurasi yang diimplementasikan pada model 
dengan parameter yang sudah di kuantitasi memiliki hasil akurasi yang sama seperti 
akurasi model yang diimplementasikan sebelum dikuantitasi.

Sedangkan pada penelitian yang dilakukan oleh Qiu et al. implemetasi CNN pada FPGA dilakukan
kuantitasi 4-bit dan 8-bit mengalami peneurunan akurasi sebesar $0.4\%$ pada model VGG16\cite{Qiu2016GoingDeeperFPGA}. 
Sedangkan pada implementasi model \emph{VGG16-SVD} yang menghasilkan akurasi sebesar $86.66\%$ dengan 
mengkuantitasi parameter 16-bit. Penelitian ini melakukan prediksi pada data citra dari \emph{image net}. 

Implementasi CNN dilakukan pada arsitektur model \emph{AlexNet}, \emph{ResNet-50}, \emph{ResNet-152}, 
\emph{RetinaNet} dan \emph{Light-weight RetinaNet} pada FPGA memakan waktu dalam skala milisekon. 
Penelitian ini dilakukan oleh Dua et al. dengan menggunakan resource DSP secara penuh pada board FPGA 
\cite{Dua2020SystolicCNN}. Hasil penelitian ini juga menunjukan bahwa \emph{convolver} bersifat \emph{compute centric} 
sedangkan \emph{neural network} bersifat \emph{memory centric}. Hal ini membuat proses konvolusi memerlukan DSP yang lebih 
banyak dibandingkan dengan \emph{neural network}. Sedangkan \emph{neural network} lebih banyak menggunakan BRAM yang 
lebih banyak.

Intrgrasi FPGA untuk melakukan \emph{stream} dalam penelitian yang dilakukan oleh Navaneethan et al. \cite{Navaneethan2022OV7670VGA} 
menggunakan \emph{VGA} dan kamera \emph{OV7670} pada board \emph{Spartan-6} dan \emph{ARTIX-7}. 
Proses vidio \emph{stream} dengan menggunakan \emph{Xilinx Vivado} dan \emph{verilog} dengan menggunakan 530 \emph{slice registers}. 
Selain itu pada penelitian Solovyev et al. juga melakukan kamera yang sama dalam melakukan prediksi inferensi pada board FPGA \emph{DEO Nano}
\cite{Solovyev2019FixedPointCNN}. Kedua penelitian ini melakukan \emph{frame buffer} pada DDR board FPGA sebelum ditampilkan ke layar. 
Alasan utama hal penggunaan DDR ini adalah ukuran \emph{frame} tidak bisa disimpan pada BRAM board FPGA terutama jika mengintegrasikan 
CNN.

Penggunaan AXI pada penelitian yang dilakukan oleh Aqueel et al. digunakan untuk 
melakukan akses menulis dan membaca memory DDR3 lewat MIG \cite{Aqueel2011DDR3Controller}.
AXI yang digunakan adalah AXI yang mendukung transaksi \emph{burst} data. 
AXI dan MIG berjalan pada sistem clock yang sama. Sistem dapat berjalan setelah MIG dapat 
terkalibrasi. AXI bertindak sebagai \emph{master} terhadap MIG sebagai \emph{slave}. 

Konfigurasi kamera dilakukan melalui \emph{Serial Camera Control Bus} atau \emph{SCCB}, yaitu 
antarmuka komunikasi yang digunakan untuk mengatur parameter internal kamera seperti format 
keluaran citra, resolusi, pengaturan warna, serta mode operasi sensor \cite{OmniVision2006OV7670}. 
Selain jalur konfigurasi, proses pembacaan data citra dari \emph{OV7670} dilakukan 
melalui sinyal \emph{pixel clock}, data paralel, serta sinyal sinkronisasi seperti \emph{HREF} dan 
\emph{VSYNC}. Sinyal \emph{HREF} digunakan untuk menandai bahwa data piksel pada satu baris sedang 
valid, sedangkan sinyal \emph{VSYNC} digunakan sebagai penanda pergantian \emph{frame}.

Tampilan \emph{VGA} dibentuk melalui pengiriman data warna RGB yang disinkronkan dengan sinyal 
\emph{horizontal sync} dan \emph{vertical sync}. Sinyal \emph{horizontal sync} berfungsi 
mengatur perpindahan baris tampilan, sedangkan \emph{vertical sync} berfungsi mengatur 
perpindahan antar-\emph{frame} pada layar \cite{Dodi2022FPGARealTimeCameraVGA}. Dalam implementasi FPGA, modul \emph{VGA controller} 
perlu menghasilkan \emph{timing} yang sesuai dengan resolusi tampilan yang digunakan, kemudian 
membaca data piksel dari \emph{frame buffer} untuk dikirimkan ke monitor.


\section{Landasan Teori}
Ne per tota mollis suscipit. Ullum labitur vim ut, ea dicit eleifend dissentias sit. Duis praesent expetenda ne sed. Sit et labitur albucius elaboraret. Ceteros efficiantur mei ad. Hendrerit vulputate democritum est at, quem veniam ne has, mea te malis ignota volumus.

Eros reprimique vim no. Alii legendos volutpat in sed, sit enim nemore labores no. No odio decore causae has. Vim te falli libris neglegentur, eam in tempor delectus dignissim, nam hinc dictas an.

Pro omnium incorrupte ea. Elitr eirmod ei qui, ex partem causae disputationi nec. Amet dicant no vis, eum modo omnes quaeque ad, antiopam evertitur reprehendunt pro ut. Nulla inermis est ne. Choro insolens mel ne, eos labitur nusquam eu, nec deserunt reformidans ut. His etiam copiosae principes te, sit brute atqui definiebas id.

  \subsection{Dasar Teori}
  Pro omnium incorrupte ea. Elitr eirmod ei qui, ex partem causae disputationi nec. Amet dicant no vis, eum modo omnes quaeque ad, antiopam evertitur reprehendunt pro ut. Nulla inermis est ne. Choro insolens mel ne, eos labitur nusquam eu, nec deserunt reformidans ut. His etiam copiosae principes te, sit brute atqui definiebas id.

  \subsection{Analisis Perbandingan Metode}
  Pro omnium incorrupte ea. Elitr eirmod ei qui, ex partem causae disputationi nec. Amet dicant no vis, eum modo omnes quaeque ad, antiopam evertitur reprehendunt pro ut. Nulla inermis est ne. Choro insolens mel ne, eos labitur nusquam eu, nec deserunt reformidans ut. His etiam copiosae principes te, sit brute atqui definiebas id.